% Шрифты и языковые настройки для XeLaTeX.
%
% Под XeLaTeX UTF-8 — родная кодировка, а Cyrillic-поддержка обеспечена
% через polyglossia + fontspec. Это избавляет от старой связки
% inputenc + fontenc(T1,T2A) + babel + cmap + textgreek + newunicodechar
% и заодно решает проблему box-drawing символов в логах TIM (одна
% строчка monofont вместо ручных \newunicodechar).

\usepackage{fontspec}
\usepackage{polyglossia}

\setdefaultlanguage{russian}
\setotherlanguage{english}

% Шрифты: оставляем CMU-семейство (визуально близкое к привычному
% Computer Modern), но monospace берём DejaVu Sans Mono — у CMU
% Typewriter в Cyrillic-наборе нет полного покрытия box-drawing
% символов, которые встречаются в выводе TIM (┌ ─ ┐ │ └ ┘ и т.п.).
\setmainfont{CMU Serif}
\setsansfont{CMU Sans Serif}
\setmonofont{DejaVu Sans Mono}[Scale=MatchLowercase]

% Тело документа набираем sans-serif (как в исходной pdflatex-сборке).
\renewcommand{\familydefault}{\sfdefault}

\usepackage{xstring}
\usepackage[autostyle]{csquotes}
\usepackage[hang,flushmargin]{footmisc}

% Меньший кегль сносок.
\renewcommand{\footnotesize}{\fontsize{6}{7}\selectfont}
